\chapter{Source Code Listings}
\label{app:code}

The complete \caesar{} source code is provided as supplementary
material and is organised as follows:

\begin{description}
  \item[\texttt{caesar/environment.py}]
    \texttt{CyberEnvironment} class: network simulation, attack injection,
    defense deployment, state vectorisation.

  \item[\texttt{caesar/ta\_gan.py}]
    \tagan{}: \texttt{ThreatAwareGenerator}, \texttt{ThreatAwareDiscriminator},
    \texttt{ThreatAwareGAN} training and inference.

  \item[\texttt{caesar/adpn.py}]
    \adpn{}: \texttt{DuelingDQN}, \texttt{ReplayBuffer}, \texttt{ADPN} agent.

  \item[\texttt{caesar/threat\_graph.py}]
    \tagraph{}: directed graph construction, Markov prediction,
    proactive defense, PageRank, effectiveness matrix.

  \item[\texttt{caesar/caesar\_algorithm.py}]
    Core co-evolutionary loop: \texttt{CAESAR.train()},
    \texttt{CAESAR.evaluate()}, fitness scoring.

  \item[\texttt{caesar/diffusion\_module.py}]
    \mgae{}: cosine schedule, score network, DDIM reverse,
    \texttt{MGAEEngine.perturb()}, robustness evaluator.

  \item[\texttt{caesar/self\_healing.py}]
    \arl{}: \texttt{SelfHealingSystem} state machine,
    countermeasure rotation, healing confirmation.

  \item[\texttt{caesar\_demo.py}]
    Complete NumPy standalone implementation for CPU-only execution.

  \item[\texttt{main.py}]
    PyTorch full-pipeline entry point with CLI arguments.

  \item[\texttt{colab/CAESAR\_Complete.ipynb}]
    Google Colab notebook with GPU acceleration and Kaggle integration.
\end{description}

All modules are documented with docstrings following the NumPy
documentation standard. Unit tests are provided in \texttt{tests/}.
